\section{Kvantinių skaičiavimų pagrindai}
\label{chap:kvantiniu-skaiciavimu-pagrindai}

Šiame skyriuje pateikiami tik tie kvantinių skaičiavimų pagrindai, kurie reikalingi suprasti darbe naudojamam Groverio algoritmo pavyzdžiui ir QPU instrukcijų paskirčiai. Nesiekiama išsamiai nagrinėti kvantinės mechanikos ar fizinių kvantinių kompiuterių realizacijų; šiame darbe kvantinis registras yra programuotojui matoma abstrakcija, kuri emuliatoriuje realizuojama per kvantinių būsenų simuliavimo biblioteką.

Groverio algoritmui suprasti būtini keli konceptai: kubitas, kelių kubitų būsenos vektorius, amplitudės ir jų tikimybės, unitarių vartų veikimas, matavimas, oracle funkcija ir amplitudžių stiprinimas. Šie konceptai yra standartinė kvantinių skaičiavimų įžanga \cite{nielsen-chuang-2010,mermin-2007}, tačiau šiame darbe jie pateikiami iš sistemos programuotojo perspektyvos: kaip operacijos, kurias turi būti įmanoma iškviesti iš RISC-V programos per QPU instrukcijas.

\subsection{Kubitai ir kvantinis registras}
\label{sec:kubitai-ir-kvantinis-registras}

Klasikinis bitas gali turėti vieną iš dviejų reikšmių: 0 arba 1. Kubitas taip pat turi dvi bazines būsenas, žymimas $\lvert 0 \rangle$ ir $\lvert 1 \rangle$, tačiau bendra kubito būsena aprašoma šių bazinių būsenų tiesine kombinacija:

\begin{equation}
    \lvert \psi \rangle = \alpha \lvert 0 \rangle + \beta \lvert 1 \rangle.
\end{equation}

Čia $\alpha$ ir $\beta$ yra kompleksiniai skaičiai, vadinami amplitudėmis. Tikimybės gaunamos iš amplitudžių modulių kvadratų, todėl galioja normalizacijos sąlyga:

\begin{equation}
    \lvert \alpha \rvert^2 + \lvert \beta \rvert^2 = 1.
\end{equation}

Kelių kubitų registras aprašomas visų bazinių būsenų amplitudėmis. Jeigu registre yra $n$ kubitų, galimų bazinių būsenų skaičius yra $2^n$, todėl bendra būsena gali būti užrašoma taip:

\begin{equation}
    \lvert \psi \rangle = \sum_{x=0}^{2^n-1} \alpha_x \lvert x \rangle.
\end{equation}

Ši eksponentinė būsenų erdvė yra viena iš priežasčių, kodėl kvantinių algoritmų simuliavimas klasikinėje mašinoje greitai brangsta. Emuliatoriuje QPU nėra fizinis kvantinis įrenginys; jis naudoja \texttt{libquantum} bibliotekos \texttt{quantum\_reg} struktūras. Todėl QPU operacijos šiame darbe reiškia ne realaus kvantinio procesoriaus valdymą, o kvantinio registro būsenos simuliavimą host sistemoje.

Groverio algoritme registras naudojamas paieškos erdvei aprašyti. Jeigu ieškoma reikšmė telpa į $n$ bitų, paieškos erdvės dydis yra $N = 2^n$. Algoritmo tikslas -- padidinti ieškomos reikšmės amplitudę taip, kad matavimo metu ji būtų gauta su didele tikimybe \cite{grover-1996}.

\subsection{Kvantiniai vartai}
\label{sec:kvantiniai-vartai}

Kvantiniai algoritmai sudaromi iš vartų, kurie keičia kvantinio registro amplitudes. Teoriškai kvantiniai vartai aprašomi unitariomis matricomis, tačiau šiame darbe svarbiausia jų programinė semantika: kiekvienas vartas yra operacija, kurią galima pritaikyti vienam ar keliems kubitams.

Groverio algoritmo realizacijoje naudojami šie vartai:

\begin{itemize}
    \item \textbf{Pauli-X vartai} -- veikia kaip kvantinis NOT: $\lvert 0 \rangle$ pakeičia į $\lvert 1 \rangle$, o $\lvert 1 \rangle$ pakeičia į $\lvert 0 \rangle$. Programoje jie naudojami pažymėtos būsenos paruošimui oracle dalyje.
    \item \textbf{Hadamard vartai} -- sukuria superpoziciją. Pritaikius šį vartą $\lvert 0 \rangle$ būsenai, gaunama vienoda $\lvert 0 \rangle$ ir $\lvert 1 \rangle$ amplitudžių kombinacija:
    \begin{equation}
        H\lvert 0 \rangle = \frac{\lvert 0 \rangle + \lvert 1 \rangle}{\sqrt{2}}.
    \end{equation}
    Pritaikius Hadamard vartus visiems paieškos registro kubitams, gaunama vienoda visų $2^n$ būsenų superpozicija.
    \item \textbf{CNOT vartai} -- dviejų kubitų vartai, kuriuose vienas kubitas yra kontrolinis, o kitas -- taikinys. Taikinys apverčiamas tik tada, kai kontrolinis kubitas yra būsenoje $1$.
    \item \textbf{Toffoli vartai} -- trijų kubitų vartai, dar vadinami kontroliuojamu kontroliuojamu NOT. Taikinys apverčiamas tik tada, kai abu kontroliniai kubitai yra $1$. Šis vartas naudojamas kelių sąlygų oracle ir difuzijos operatoriaus konstrukcijai.
\end{itemize}

\subsection{Matavimas ir tikimybės}
\label{sec:matavimas-ir-tikimybes}

Kvantinio registro būsena nėra tiesiogiai perskaitoma kaip klasikinis atminties masyvas. Matavimas grąžina vieną iš bazinių būsenų, o konkretaus rezultato tikimybė priklauso nuo tos būsenos amplitudės. Jeigu bazinės būsenos $\lvert x \rangle$ amplitudė yra $\alpha_x$, tikimybė gauti $x$ matavimo metu yra:

\begin{equation}
    P(x) = \lvert \alpha_x \rvert^2.
\end{equation}

Kadangi amplitudės yra kompleksiniai skaičiai, praktinėje realizacijoje jos dažnai saugomos kaip realioji ir menamoji dalys. Jeigu $\alpha_x = a + bi$, tikimybė apskaičiuojama taip:

\begin{equation}
    P(x) = a^2 + b^2.
\end{equation}

Darbe realizuotame QPU plėtinyje šis ryšys matomas per kelias instrukcijas. \texttt{Q.GETREGNODE} leidžia nuskaityti \texttt{libquantum} registro mazgą: bazinę būseną, realią amplitudės dalį ir menamą amplitudės dalį. \texttt{Q.PROB} iš realiosios ir menamosios dalių apskaičiuoja tikimybę. \texttt{Q.BMEASURE} atlieka vieno kubito matavimą ir pakeičia registro būseną pagal matavimo rezultatą.

Groverio algoritmo pavyzdyje po iteracijų nuskaitomi registro mazgai ir ieškoma tokios būsenos, kuri sutampa su ieškoma reikšme. Kai tokia būsena randama, jos amplitudė paverčiama tikimybe ir išvedama vartotojui. Todėl eksperimento rezultatas nėra vien programos baigtis; svarbu ir tai, kokia tikimybė priskiriama ieškomai reikšmei.

\subsection{Groverio algoritmo idėja}
\label{sec:groverio-algoritmo-ideja}

Groverio algoritmas sprendžia nestruktūrizuotos paieškos problemą: duotoje $N$ elementų erdvėje reikia rasti pažymėtą elementą. Klasikinėje paieškoje blogiausiu atveju reikia patikrinti $O(N)$ elementų, o Groverio algoritmas pažymėtą elementą randa naudojant apie $O(\sqrt{N})$ oracle užklausų \cite{grover-1996,nielsen-chuang-2010}.

Algoritmas prasideda nuo vienodos superpozicijos paruošimo. Jeigu paieškos registre yra $n$ kubitų, po Hadamard vartų visoms paieškos kubitų pozicijoms gaunama būsena:

\begin{equation}
    \frac{1}{\sqrt{N}} \sum_{x=0}^{N-1} \lvert x \rangle,
    \quad N = 2^n.
\end{equation}

Tada kartojami du pagrindiniai veiksmai:

\begin{itemize}
    \item \textbf{oracle} -- operacija, kuri pažymi ieškomą būseną. Dažnai tai aprašoma kaip pažymėtos būsenos fazės pakeitimas, t. y. ieškomos būsenos amplitudė dauginama iš $-1$;
    \item \textbf{difuzijos operatorius} -- operacija, vadinama amplitudžių inversija apie vidurkį. Ji padidina pažymėtos būsenos amplitudę ir sumažina kitų būsenų amplitudes.
\end{itemize}

Vienas oracle ir difuzijos operatoriaus pritaikymas sudaro vieną Groverio iteraciją. Kai pažymėtas vienas elementas, apytikslis iteracijų skaičius yra:

\begin{equation}
    k \approx \frac{\pi}{4}\sqrt{N}.
\end{equation}

Po iteracijų registras matuojamas arba analizuojamos jo amplitudės. Idealiu atveju pažymėtos būsenos tikimybė tampa didelė, todėl matavimas dažniausiai grąžina ieškomą elementą. Ši savybė daro Groverio algoritmą tinkamą šiam darbui: jis pakankamai mažas, kad būtų realizuojamas kaip vartotojo erdvės programa, bet kartu naudoja kelis svarbius kvantinius veiksmus -- superpoziciją, oracle, daugkartinę iteraciją ir tikimybių analizę.

